\subsection{Riesgos}

El Cuadro~\ref{tab:gestion_riesgos} reúne riesgos de diseño y medidas para mitigarlos; no informa mediciones ni resultados de mitigación.

\begin{table}[H]
\caption{Riesgos prioritarios y medidas de mitigación}
\label{tab:gestion_riesgos}
\begin{tabularx}{\textwidth}{|X|X|}
\hline
\textbf{Riesgo} & \textbf{Medida de mitigación} \\
\hline
Transcripción incorrecta o ambigua. & Permitir revisión y corrección; evaluar ASR local con datos sintéticos y dispositivos representativos. \\
\hline
Acceso clínico no autorizado. & Aplicar permisos efectivos, institución y asignación asistencial mediante RLS; verificar respuestas denegadas. \\
\hline
Discrepancia entre controles de interfaz y autorización efectiva. & Mantener RLS como autoridad de acceso y probar permisos, institución y asignación en backend. \\
\hline
Exposición de audio o transcripción no confirmada. & Mantenerlos en el dispositivo, eliminar temporales conforme a política y comprobar el flujo de red y ciclo de vida. \\
\hline
Interrupción de consulta o persistencia. & Comunicar fallos y evitar pérdida o duplicación silenciosa; la mitigación no presupone operación offline completa. \\
\hline
Indisponibilidad de los servicios de backend o de la fuente. & Comunicar el estado y preservar la revisión y autorización; evaluar recuperación según el entorno. \\
\hline
Propuesta de redacción o resumen incorrecta. & Ejecutar asistencia local, limitarla a fuentes autorizadas y someter el contenido a revisión humana; no publicar automáticamente. \\
\hline
Consulta o vinculación incorrecta de información externa. & Validar los contratos de las fuentes y la vinculación con hospitalización mediante pruebas con datos sintéticos. \\
\hline
Baja comprensión de los recorridos. & Aplicar el protocolo definido si se realiza una evaluación; documentar método y observaciones sin presuponer participación o resultado. \\
\hline
\end{tabularx}
\end{table}

La infraestructura administrada no elimina por sí sola los riesgos. Autorización, privacidad, revisión humana y ciclo de vida requieren controles en Flutter, la pasarela y Supabase; su eficacia debe sustentarse con evidencia.
